% !TEX program = pdflatex
%
% Proof Without a Verdict — an ACTIK case
% Preprint. Comments welcome.
%
% Build:  make -C paper        (single column and two column)
% Two column only:  make -C paper twocol
\ifdefined\twocol
  \documentclass[10pt,a4paper,twocolumn]{article}
\else
  \documentclass[11pt,a4paper]{article}
\fi

\usepackage[T1]{fontenc}
\usepackage{lmodern}
\usepackage[utf8]{inputenc}
\ifdefined\twocol
  \usepackage[margin=0.85in]{geometry}
  \setlength{\columnsep}{20pt}
\else
  \usepackage[margin=1in]{geometry}
\fi
\usepackage[protrusion=true,expansion=false]{microtype}
\usepackage{booktabs}
\usepackage{tabularx}
\usepackage{xcolor}
\usepackage[numbers,sort&compress]{natbib}
\usepackage{enumitem}
\usepackage[hyphens]{url}
\usepackage[hidelinks,breaklinks]{hyperref}
\providecommand{\doi}[1]{\href{https://doi.org/#1}{doi:\nolinkurl{#1}}}

\makeatletter
\g@addto@macro\UrlBreaks{%
  \do\a\do\b\do\c\do\d\do\e\do\f\do\g\do\h\do\i\do\j\do\k\do\l\do\m%
  \do\n\do\o\do\p\do\q\do\r\do\s\do\t\do\u\do\v\do\w\do\x\do\y\do\z%
  \do\0\do\1\do\2\do\3\do\4\do\5\do\6\do\7\do\8\do\9\do\-\do\.\do\=\do\&\do\?}
\makeatother
\Urlmuskip=0mu plus 1mu

\hypersetup{
  pdftitle={Proof Without a Verdict},
  pdfauthor={Sengtha Chay; Mouniseyhak Va},
  pdfsubject={Holder-controlled credentials, limits on what may be asked, and tiered issuers: an ACTIK case},
}

\newcommand{\code}[1]{\texttt{\small #1}}
\newcommand{\kind}[1]{\textsc{#1}}

\ifdefined\twocol
  \newenvironment{widetable}{\begin{table*}[t]}{\end{table*}}
  \setlength{\emergencystretch}{3em}
  \tolerance=1200
\else
  \newenvironment{widetable}{\begin{table}[t]}{\end{table}}
\fi

% ---------------------------------------------------------------------------
% AUTHORS: the author list, affiliations and the contribution statement below
% must be confirmed by the people named before this is posted.
% ---------------------------------------------------------------------------
\title{\textbf{Proof Without a Verdict}\\[0.4em]
  \large Holder-Controlled Credentials, Limits on What May Be Asked,\\
  and Tiered Issuers --- An ACTIK Case}

\author{Sengtha Chay\textsuperscript{1} \and Mouniseyhak Va\textsuperscript{2}\\[0.6em]
  \small National University of Management, Phnom Penh, Cambodia\\
  \small \textsuperscript{1}Director, CamboVerse Center \quad \textsuperscript{2}CamboVerse DevOps\\
  \small Corresponding author: \texttt{sengtha@num.edu.kh}}

\date{October 2026\\[0.5em]\normalsize\emph{Preprint --- comments welcome.}}

\begin{document}
\maketitle

\noindent\textbf{Keywords:} verifiable credentials; selective disclosure; holder binding;
credential revocation; recruitment; non-discrimination; digital certificates;
employment records; trust lists; emerging economies; CamboVerse; Cambodia

\paragraph{Disclaimer.}
This is independent academic research. The views are the authors' own and do
not represent the National University of Management or any body within it. The
paper is not affiliated with, endorsed by, funded by, commissioned by, or
developed for any government, ministry, accreditation body, regulator, employer
or other institution. The roles named in the design --- Root, accredited
institution, registered employer, holder, requester --- are hypothetical
placeholders; nothing here implies that any body holds, or has agreed to hold,
any of them. Institution names in examples are fictional unless a source is
cited. ACTIK and CamboVerse are project names. No issuance, adoption, usage or
hiring figures appear in this paper, because none exist: ACTIK has not been
deployed. What has been built has been tested, and the tests are public.

\begin{abstract}
\noindent
Systems that verify certificates usually answer with a verdict: a tick, a
green screen, the word \emph{verified}. We argue that the verdict is the wrong
output, and that the questions a verifier is asked to settle are set by the
format before anyone looks at a credential. We make the argument with ACTIK, an
open-source application for issuing and managing digital assets --- certificates
first, employment records second --- built for Cambodia and intended as the data
layer of the CamboVerse ecosystem.

Four design choices carry the argument. First, verification returns evidence
rather than a verdict: who issued the credential, its standing as of a dated
list, and the fields a reader must compare with the document in front of them;
the result types contain no boolean. Second, an employer's request for proof is
written in a format that cannot ask for date of birth, sex, marital status, a
photograph, a national identity number or a student number, and the limit is
enforced three times --- in the candidate's application, in the database, and in
the employer's application --- so that it does not depend on a registrar. Third,
the signed trust list carries each issuer's kind: a registered employer may sign
employment records and nothing else, and every verifier says which kind of
issuer it is. Fourth, withdrawal lists name no one, and a credential issued to a
person who already holds a wallet is bound to that wallet's key, so a copy or a
replayed answer is refused.

The cryptography is standard and we claim none of it. Selective disclosure, key
binding, status lists and limits on over-asking all exist. Cambodia's national
platform already verifies official documents online, against a record it holds;
what it does not do is put a credential in the holder's own hand, from any
accredited issuer, to be checked offline and shown in part. ACTIK does, and is
offered beside the national platform as a choice, not in place of it. Our
contribution is narrower than the mechanisms: where the
limits sit, what the verifier is told, and which of the claims a design makes
can be checked in its code. We separate claims that are true of the published
system, claims carried from earlier work, claims that follow from stated
premises, and guesses to be tested. Nothing has been issued to anyone.
\end{abstract}

% ===========================================================================
\section{Introduction}
\label{sec:intro}

A student finishes a course on a digital topic and asks a simple question: can I
have my certificate digitally --- not a scan, not a link to someone else's
website, but a certificate I keep, that anyone can check, and that I can show
without showing everything else? In Cambodia today the answer is mostly no. The
national verification platform checks official documents against a record held
on its own servers; it gives the holder a code that points there, not a
credential of their own~\citep{verifygovkh,chay2026qrseal}. A course from a
training provider usually ends with a PDF or a sheet of paper.

Later the same person applies for a job. She brings a photocopy of her
degree, certified as a true copy, and the employer files it. Nobody checks it
against the university, because there is no easy way to; and the employer, if
it wants to, can ask her age, whether she is married, and for a photograph,
because the application form has spaces for them.

Two things are wrong in that scene, and they are usually treated separately.
The first is that the document cannot be checked. The second is that the
employer can ask for more than it needs. Technology for the first problem is
mature: a signed credential can be checked by anyone, offline, in milliseconds.
The second problem is usually left to law and to the employer's conscience.

This paper is about a system that gives people their certificates digitally,
and about the place where those two problems meet. When a candidate answers an
employer with signed credentials instead of a form, the format of the request
decides what can be asked, and the format of the answer decides what the
employer is told. Both are design decisions, made once, and they shape every
exchange that follows.

\subsection{What already exists}

We start by conceding most of the ground. Verifiable credentials are a W3C
Recommendation~\citep{vcdm2}. Selective disclosure, in which the holder reveals
some signed fields and withholds others, is an IETF standard, SD-JWT, which also
specifies how a presentation is bound to the holder's key~\citep{rfc9901,rfc7800}.
Status lists for revoking credentials without a per-credential lookup are a W3C
Recommendation~\citep{bitstring}. The European Union's digital-identity framework
already obliges the parties that request data from a wallet to register what they
intend to ask for, and forbids them to ask for more~\citep{eidas2,eu2025848}. And
Cambodia already has a national document-verification service, Verify.gov.kh,
which places QR codes on diplomas and other official documents and checks them
online against a record held by the platform~\citep{verifygovkh,cambodianess2026verifygov,khmertimes2026verifygov}.

None of that is ours. What the national platform does not offer --- a credential
the holder keeps and presents, from issuers beyond government bodies, checked
without contacting anyone --- is what ACTIK adds, and it is offered beside the
platform, for people and institutions to choose, not instead of it
(Section~\ref{sec:choice}).

\subsection{What we claim}
\label{sec:claims}

We make four claims, and for each we say what is already known.

\textbf{First: the output should be evidence, not a verdict.} A verifier that
says \emph{verified} tells the reader that a registered key signed some bytes,
and invites them to believe much more: that the person in front of them is the
person named, that a two-week course and a four-year degree are equally
weighty, that a credential checked against last month's withdrawal list is
current today. Research on security indicators has long found that people do not
notice or act on the indicators designers rely on~\citep{schechter2007emperor}.
What we add is a concrete interface contract, enforced by the result types and
by tests: the issuer first, its kind, the standing as of a dated list, the fields
to compare, and three distinct ways of not succeeding.

\textbf{Second: limits on what may be asked belong in the format.} The EU
mechanism limits over-asking through registration: a national register records
what each requester intends to ask, and a certificate enforces
it~\citep{eu2025848}. That requires a registrar, a register and a legal
framework. We put a fixed limit in the request format itself, so that a request
for date of birth, sex, marital status, a photograph, a national identity number
or a student number cannot be written at all, and we enforce it at three points
that do not trust each other. The limit is wider than the grounds listed in
Cambodia's Labour Law~\citep{khlabourlaw}, and we say where and why
(Section~\ref{sec:known-hiring}).

\textbf{Third: the trust list should say what kind of issuer each issuer is.}
A trust list that admits an issuer usually admits it for everything. We sign each
issuer's kind into the list --- an accredited institution may issue any credential
type, a registered employer only employment records --- and every verifier shows
the kind beside the name. Employment records then carry their own rules: they
cannot hold a salary, a reason for leaving or a performance rating, and
\emph{currently employed} is always shown as of the day it was signed.

\textbf{Fourth: withdrawal should name no one, and a credential should be bound
to its holder where it can be.} Both are carried techniques. The withdrawal lists
use the hashed-entry design of an earlier printed-certificate
specification~\citep{chay2026qrseal}; binding follows SD-JWT key
binding~\citep{rfc9901}. What we add is their combination with the first three
choices, and an account of what binding does not do.

Section~\ref{sec:known} reviews what earlier work covers. Section~\ref{sec:method}
explains the method. Section~\ref{sec:system} describes the system and what has
been tested. Sections~\ref{sec:verdict}--\ref{sec:tiers} develop the four claims.
Section~\ref{sec:discussion} discusses them, Section~\ref{sec:unknown} lists what
we do not know, Section~\ref{sec:next} proposes what should happen next, and
Section~\ref{sec:conclusion} concludes.

% ===========================================================================
\section{What Is Already Known}
\label{sec:known}

\subsection{Credentials, selective disclosure and binding}

The W3C data model defines a credential as a set of claims made by an issuer
about a subject, and a presentation as data derived from credentials and shared
by a holder~\citep{vcdm2}. SD-JWT signs, instead of each claim, a salted hash of
it, so that a holder can later reveal a chosen subset; the verifier can check
each revealed claim against the signed hashes without learning the
rest~\citep{rfc9901}. The same standard defines key binding: the issuer places
the holder's public key in the credential, using the confirmation claim of
RFC~7800~\citep{rfc7800}, and each presentation carries a short token signed with
the matching private key, naming its audience and hashing the presentation it
accompanies. ACTIK uses SD-JWT and its key binding as specified. We claim nothing
about either.

\subsection{Signed documents on paper}

Signed payloads in QR codes are established: the EU Digital COVID Certificate put
a COSE-signed, compressed, base45-encoded payload on paper at continental
scale~\citep{eudcc2021,rfc9052}. The printed-certificate format ACTIK uses
(Section~\ref{sec:printed}) is KH-SQR Profile B, specified and implemented in an
earlier preprint by one of the present authors~\citep{chay2026qrseal,khsqr2026repo}.
That work also supplies the withdrawal-list entry design we carry: an entry is a
hash of the issuer and the document number, so a list can be public without
listing anyone.

\subsection{Verification in Cambodia}

Verify.gov.kh, the national platform, gives an issued document a QR code and a
digital code which, with condensed information from the document, are stored on a
blockchain; opening the code shows the stored contents, and the platform checks
that the document comes from a recognised government body~\citep{verifygovkh}.
Press reports describe its use for grade-12 certificates and university degrees,
including degrees of the university at which the present authors work, and its
adoption beyond Cambodia~\citep{khmertimes2026verifygov,cambodianess2026verifygov}.
We have not verified those figures and do not rely on them.

Its architecture has been analysed in an earlier preprint by one of the present
authors, and we carry that analysis rather than repeat it~\citep{chay2026qrseal}.
In short: the code on an issued document carries a \emph{reference} to a record,
a URL with a key in it, rather than the credential itself. Verification therefore
depends on the platform staying reachable for the life of the document; a
photograph of the document carries a working key; and official codes that open
websites teach the public a habit that QR phishing exploits. The same analysis
credits what a reference buys --- a correction or withdrawal is visible to the
next reader at once --- and we return to that in Section~\ref{sec:choice}.

\subsection{Limits on what a requester may ask}
\label{sec:known-hiring}

Two bodies of rule bear on what an employer may ask a candidate. The first is
data protection. Under the EU framework, a party that requests attributes from a
digital-identity wallet must register its intended use in a national register,
must not request more than it registered, and is identified to the wallet by a
registration certificate~\citep{eidas2,eu2025848}. This is a limit by
registration: flexible, per requester, and dependent on a registrar.

The second is labour law. Cambodia's Labour Law, in its Article~12, prohibits an
employer from taking account of race, colour, sex, creed, religion, political
opinion, birth, social origin, or trade-union membership and activity in hiring
and later employment decisions~\citep{khlabourlaw}. We could not reach the
promulgated text from the environment in which this paper was prepared, and
quote the list from secondary legal summaries; it should be checked before
reliance. On that list, age and marital status do not appear. ACTIK's request
format excludes them anyway, together with a photograph, a national identity
number, a student number, place of birth and contact details. That is a design
choice, not a legal requirement, and Section~\ref{sec:asking} argues for it.

% ===========================================================================
\section{Method}
\label{sec:method}

\subsection{What kind of paper this is}

This is a design case with a conceptual argument. In Jaakkola's terms it is
closest to a theory adaptation: it takes ideas established for other purposes ---
credential formats, security-indicator research, data-protection limits on
requesters --- and applies them to a setting they were not built
for~\citep{jaakkola2020}. It follows the method of a companion CamboVerse case by
several of the present authors~\citep{chay2026precinct}, including its four kinds
of claim (Section~\ref{sec:fourkinds}).

\subsection{How we use the case}

Our case is ACTIK, an application whose purpose is to issue digital assets and to
let their owners manage them. A digital certificate is the first kind of asset it
handles; employment records are the second. ACTIK is planned as the data layer of
CamboVerse, an open-source Virtual Cambodia: under the agreed division, ACTIK holds
identities, issued files, verification state and withdrawal, and CamboVerse
displays what a person chooses to show~\citep{chay2026precinct}.

We use ACTIK as an illustration, not as a test. Every design claim in
Section~\ref{sec:system} corresponds to code a reader can inspect and, in most
cases, to a test that fails if the claim stops being true~\citep{actik2026repo}.
Nothing has been issued to a real person, no employer has used it, and no
statement here is about how people behave with it.

\subsection{What we checked before claiming anything}

Before settling the claims we searched the standards and regulation named in
Section~\ref{sec:known}. Three claims we had intended to make were
withdrawn as a result. Key binding was not ours to claim: it is in
SD-JWT~\citep{rfc9901}. Privacy-preserving revocation was not ours: status lists
address it~\citep{bitstring}, and our entry design is carried from the
printed-certificate work~\citep{chay2026qrseal}. And limiting what a requester
may ask was not ours either: the EU framework does it by
registration~\citep{eu2025848}. What survives is the narrower claim stated in
Section~\ref{sec:claims}.

\subsection{Four kinds of claim}
\label{sec:fourkinds}

Every substantive statement in this paper is one of four
kinds~\citep{chay2026precinct}. A \kind{specified} claim is true of the published
code and can be checked against it; where a test enforces it, we say so. A
\kind{carried} claim is established in earlier work and imported. An
\kind{analytic} claim follows from stated premises; it holds if they do. A
\kind{conjectural} claim is a reasoned guess about how people will behave that
nothing here establishes.

The line most easily blurred to our advantage is between analytic and
conjectural. That a request format which has no field for marital status cannot
carry a request for it is analytic. That employers who cannot ask for marital
status will therefore discriminate less is conjectural --- they can still ask at
interview --- and we have asked no employer.

\subsection{Who we are}

We are not independent of the case. One author directs the CamboVerse Center, an
institute at the National University of Management; the other is an engineer on
the CamboVerse project and, according to the repository history, wrote most of
ACTIK's original application. Parts of the system described in
Section~\ref{sec:system} were written with the assistance of an AI model, as the
declarations record. The university at which we work issues degrees, and its
degrees are verified through the national platform discussed in
Section~\ref{sec:known}~\citep{khmertimes2026verifygov}. These are affiliations
and facts, not claims about who governs anything.

This cuts both ways. It strengthens Section~\ref{sec:system}: every statement
there can be checked against public code and reproduced with public tests. It
weakens Sections~\ref{sec:verdict}--\ref{sec:discussion}: we have an interest in
the design looking important. The conjectural claims are where that interest
could distort our judgement most, and they are exactly the ones we have not
tested.

% ===========================================================================
\section{The System, Described}
\label{sec:system}

This section describes. It does not argue. Every statement in it is
\kind{specified}.

\subsection{Roles}

\begin{description}[leftmargin=0pt,itemsep=2pt]
\item[Root.] Holds an offline signing key whose public half is built into the
application. It signs the trust list. Approving an issuer in the database does
not make it trusted; the Root's next signed list does.
\item[Accredited institution.] An issuer on the trust list that may issue any
credential type.
\item[Registered employer.] An issuer on the trust list that may issue
employment records only.
\item[Holder.] A person with a wallet. The wallet is end-to-end encrypted: the
database holds ciphertext, and only the holder's PIN or passkey opens it.
\item[Requester.] Any signed-in account that asks holders for proof --- in
practice, an employer recruiting.
\end{description}

\subsection{Issuing and accepting}

An issuer signs an SD-JWT credential in its own browser with a key held in memory
as a non-extractable Web Crypto key~\citep{webcrypto}; the key never leaves that
object. The credential goes to an outbox addressed to the recipient's email. The
database admits it only if the signed-in user owns an issuer of the right kind
for that credential type. Nothing reaches the recipient's wallet until they
accept it, and they may decline, which deletes the offer. On acceptance the
holder's application verifies the credential against the trust list and the
issuer's withdrawal list, and stores it encrypted.

\subsection{The trust list and issuer kinds}

The trust list is a statement signed by the Root, valid for at most 31 days,
listing each issuer's decentralised identifier, name, kind and keys, with each
key's validity window and status. A verifier refuses a list signed by any other
key, an expired list, or a list older than one it has already seen. A key that
has been retired continues to vouch for what it signed before retirement; a
revoked key vouches for nothing. A registered employer that signs any credential
other than an employment record is refused with the reason
\code{TYPE\_NOT\_ALLOWED\_FOR\_ISSUER}, in the application and on printed codes.
The build step that the Root runs prints each issuer's kind for review and
refuses to sign a kind it does not recognise.

\subsection{Withdrawal}

Each issuer signs its own withdrawal list with a key the trust list accepts for
it. An entry is the SHA-256 hash of a domain tag, the issuer and either the
document number or the credential's identifier, with one of two fixed reasons
--- \emph{withdrawn} or \emph{replaced by a corrected credential} --- and a date.
The database refuses a list with any other field, so a document number or a
free-text explanation cannot be published by mistake. A verifier reports a
credential's standing as \emph{clear} against a named list version and date,
\emph{unchecked} if the issuer publishes no list or its list has lapsed, or
\emph{withdrawn}, with the reason.

\subsection{Holder binding}

The first time a wallet is unlocked, the application makes a holder key pair. The
private half is encrypted with the wallet's own key before it is stored, and is
held in memory as a non-extractable key. When a recipient already has a wallet,
the issuer writes the holder's public key into the credential as a confirmation
claim~\citep{rfc7800}. Every presentation of a bound credential must then carry a
key-binding token signed by that key, naming its audience --- one share link or
one proof request --- and hashing the presentation~\citep{rfc9901}. A verifier
refuses a bound credential presented without the token, with a token from any
other key, with a token dated in the future or not matching the disclosures, or
with a token made for another audience. A credential issued before the recipient
had a wallet is unbound, still verifies, and is labelled as unbound.

\subsection{Sharing and verifying}

A holder can share a credential through a link. The share discloses a fixed set
of fields --- the holder's name, the issuer, the document number and date, and the
substance of the credential --- plus any optional fields the holder selects, and
can be limited to a single view, enforced by the database. The verifier's page
checks the presentation on the verifier's device and shows the issuer first, its
kind, the standing, the four fields to compare with any paper document, and
whether the presentation was bound.

\subsection{Proof requests}
\label{sec:proofrequests}

A requester writes a request naming up to five credential types, each with
optional extra fields from a short list for that type, and an expiry of at most
90 days, and shares its link. Table~\ref{tab:allowlist} gives the list. A request
cannot name any field outside it. A candidate opens the link, signs in, and sees
the matching credentials in their wallet, each already checked; a withdrawn one
cannot be chosen. The candidate sees exactly what the requester will see before
sending, and chooses the contact details they leave. The answer can be withdrawn
by the candidate, never edited, and is visible only to the candidate and the
requester. Nobody can list requests or browse candidates.

\begin{widetable}
\centering\small
\caption{What a proof request can ask for. Every answer discloses the fields in
the middle column; a request may add those in the right column. Nothing else can
be requested. (\kind{specified}; enforced by \code{test-proof.ts} and by the
database guard.)}
\label{tab:allowlist}
\begin{tabularx}{\linewidth}{@{}lXl@{}}
\toprule
Credential type & Always disclosed & May also be asked \\
\midrule
Academic degree & name, institution, degree, graduation date, certificate number & major, GPA \\
Professional certification & name, institution, certification, issuing body, date, licence number, expiry & --- \\
Certificate of completion & name, institution, type, programme, completion date & duration, department or role \\
Attendance or participation & name, institution, type, event, date, organiser & role \\
Merit or excellence award & name, institution, achievement, date awarded & basis of award \\
Appreciation or service & name, institution, type, reason, date & capacity \\
Employment record & name, employer, job title, employment type, start, status, end & department, role \\
\midrule
\multicolumn{3}{@{}p{\linewidth}@{}}{\emph{Never requestable:} date of birth or age, sex or
gender, marital status, photograph, national identity number, student number,
place of birth, religion or ethnicity, email or phone number.}\\
\bottomrule
\end{tabularx}
\end{widetable}

The limit is enforced three times. The candidate's application discloses only
allowed fields, whatever the request says. The database decodes every disclosure
in every answer and refuses an answer that carries any other claim, or a
credential of the wrong type. The requester's application refuses to display an
answer that discloses anything else --- not even its allowed part.

\subsection{Employment records}

An employment record carries the holder's name, the employer, a job title, an
employment type, a start date and a status --- \emph{current} or \emph{ended},
with an end date --- and optionally a department and a one-line role. The issuing
form has no field for salary, reason for leaving, performance or discipline, and
no free-text note or attached file; the database decodes each record and refuses
any claim outside that set. \emph{Current} is displayed everywhere as
\emph{current, as of} the date the employer signed it. When a job ends, the
employer issues a new record and withdraws the old one as replaced.

\subsection{Printed certificates}
\label{sec:printed}

For degrees and professional certifications that carry a document number, the
issuer also signs a printed code in the KH-SQR Profile B format: CBOR claims in a
COSE signature, compressed and base45-encoded behind a \code{KH1:}
prefix~\citep{chay2026qrseal,rfc9052}. The code is delivered to the holder with
the credential, checked on acceptance, and kept in the wallet for reprinting.
ACTIK's scanner verifies it offline against the same trust and withdrawal lists.
The vendored implementation passes all thirteen Profile B verification vectors of
the original specification.

\subsection{The seam with CamboVerse}

A holder can export a credential as an exhibit for a CamboVerse personal museum.
The exhibit carries the SHA-256 of the original file, a picture the holder has
cropped and covered, reduced and re-encoded on their device, the printed code if
the holder chooses, and the withdrawal-list hashes under which the credential
would appear. It carries no account identity, no original file and no
\emph{verified} flag; its visibility defaults to private, and CamboVerse is to show
the stricter of the exhibit's and the room's setting.

\subsection{What is tested and measured}
\label{sec:tested}

Table~\ref{tab:tests} lists the automated suites. Database rules were exercised
against a local PostgreSQL~16 instance with the platform's authentication schema
stubbed, and interface behaviour in headless Chromium against the built
application with a stand-in for the hosted database; neither is the hosted
service. Table~\ref{tab:measure} reports sizes and timings produced by a script in
the repository.

\begin{widetable}
\centering\small
\caption{Automated test suites (\code{npm test}). Counts are of assertions on the
cited branch~\citep{actik2026repo}.}
\label{tab:tests}
\begin{tabularx}{\linewidth}{@{}lrX@{}}
\toprule
Suite & Assertions & What it holds to \\
\midrule
\code{test-sdjwt.ts} & 21 & issuance, selective disclosure, signature and expiry checks \\
\code{test-claim.ts} & 11 & nothing enters a wallet unless the trust list's key signed it \\
\code{test-trust.ts} & 48 & Root-signed trust list, rollback, key windows; withdrawal lists that name no one \\
\code{test-printed.ts} & 48 & the thirteen Profile B vectors; printed codes on the trust layer \\
\code{test-museum.ts} & 19 & exhibit export carries no original file and no verdict \\
\code{test-proof.ts} & 34 & what can be asked, what is disclosed, how an answer is checked \\
\code{test-employment.ts} & 21 & issuer kinds; what an employment record may say \\
\code{test-binding.ts} & 15 & key binding: missing, foreign, replayed and tampered proofs \\
\midrule
Total & 217 & \\
\bottomrule
\end{tabularx}
\end{widetable}

\begin{widetable}
\centering\small
\caption{Sizes and timings (\code{npm run measure}), with fresh keys and fictional
names. Timings are of Node.js~22 on a four-core 2.8\,GHz server processor, not of
a phone, and vary by about 0.2\,ms between runs; they are illustration, not a
claim about the devices the design targets.}
\label{tab:measure}
\begin{tabularx}{\linewidth}{@{}Xr@{}}
\toprule
Quantity & Value \\
\midrule
Degree credential, all claims (bytes) & 1{,}909 \\
Default share presentation (bytes) & 1{,}702 \\
Key-binding token added to it (bytes) & 366 \\
Trust list with 100 issuers (bytes) & 34{,}886 \\
Withdrawal list with 1{,}000 entries (bytes) & 118{,}135 \\
Printed code, Khmer holder name (characters) & 417 \\
Verify a bound share against a 1{,}000-entry list, median / 95th percentile (ms) & 1.7 / 2.4 \\
Open a 100-issuer trust list, median / 95th percentile (ms) & 1.2 / 1.7 \\
Verify a printed code, median / 95th percentile (ms) & 9.0 / 12.3 \\
Application shell cached for offline use (KiB) & 1{,}392 \\
\bottomrule
\end{tabularx}
\end{widetable}

\subsection{What the system does not show}

There is no deployment. No credential has been issued to a real person, no
employer has written a request, and no candidate has answered one. The browser
checks ran against a stand-in database, not the hosted service. No timing has
been taken on the low-cost phones the design targets. The Root role has been
exercised only in a rehearsal of its key ceremony; nobody holds it.

% ===========================================================================
\section{Evidence, Not a Verdict}
\label{sec:verdict}

\subsection{What a verdict hides}

A signature check establishes one fact: a key the verifier trusts signed these
bytes. Everything else a reader wants to know lies outside it. The paper document
may be a forgery carrying a genuine code copied from another certificate. The
person presenting may not be the person named. The issuer may have withdrawn the
credential after the verifier's last list. The issuer may be a three-person
company whose records weigh little, or a university. A single \emph{verified}
mark collapses all of these into one bit.

\textbf{Proposition 1} (\kind{analytic}). \emph{If a verification interface
reports a single success state, then a credential copied onto a forged document,
a credential presented by someone other than its holder, a credential withdrawn
after the last list, and a credential from a low-weight issuer are reported
identically to a credential with none of those defects.} It follows from the
premise that the signature check is the same in all five cases.

\subsection{What ACTIK reports instead}

ACTIK's result types make the verdict hard to write. A successful check returns a
record, not a flag: the issuer and its kind from the signed list, the standing
with its list version and date, the four fields a reader must compare with any
document (name, document number, issuing organisation, date), and whether the
presentation was bound to its holder (\kind{specified}). A failure is one of three
distinct outcomes: \emph{refused}, with a stable reason; \emph{could not check},
when the verifier's own trust state is at fault and the credential is not; and
\emph{withdrawn}, naming the issuer's decision (\kind{specified}). Security
indicator research suggests that a bare success mark is not read the way
designers hope~\citep{schechter2007emperor} (\kind{carried}).

\textbf{Proposition 2} (\kind{conjectural}). \emph{Readers shown the issuer,
its kind and the dated standing first will distinguish a short course from a
degree, and an unchecked standing from a clear one, more often than readers shown
a single success mark.} Nothing here tests it.

% ===========================================================================
\section{Asking Less, by Design}
\label{sec:asking}

\subsection{Limits by registration and limits by format}

A limit on over-asking can live in two places. It can live in a register: each
requester declares what it will ask, and the wallet refuses more. That is the EU
mechanism~\citep{eu2025848} (\kind{carried}). Or it can live in the format: the
request language has no way to express the forbidden field. ACTIK does the second
(\kind{specified}).

\textbf{Proposition 3} (\kind{analytic}). \emph{A limit written into the request
format binds every requester without a registrar, at the cost of being the same
for every requester.} A register can allow a bank to ask what an employer cannot.
A format cannot. We accept that cost because the only requester ACTIK is built
for is an employer judging qualifications, and because Cambodia has no register
of requesters to rely on.

\subsection{Why three enforcement points}

A limit enforced only in the candidate's application is defeated by a modified
application; one enforced only in the database is defeated by a misconfigured
database; one enforced only in the requester's application is defeated by
reading the database directly. ACTIK enforces the allowlist at all three
(\kind{specified}). This is an instance of a more general observation made in an
earlier Cambodia case by several of the present authors: a control enforced in
one layer does not travel to claims made in another~\citep{chay2026stages}
(\kind{carried}, self-cited).

\subsection{Why wider than the law}

The excluded fields are wider than the grounds in Article~12 of the Labour
Law~\citep{khlabourlaw}: age, marital status and a photograph are not on that list
as we have it. We exclude them because each is routinely asked for in job
advertisements, none bears on whether a credential is genuine, and a photograph
discloses sex, apparent age and ethnicity at once. The exclusion of a student
number and a national identity number has a different reason: they are
identifiers that join records across systems, and a qualification can be judged
without them.

\textbf{Proposition 4} (\kind{conjectural}). \emph{Employers who receive
candidates' qualifications as signed answers without age, sex, marital status or
a photograph will shortlist on qualifications more than employers who receive
paper applications.} The opposing conjecture is as reasonable: that employers
will ask the excluded questions at interview, and that the format changes the
order in which bias acts, not whether it acts. The name, which the format must
disclose so that identity can be checked, itself carries signals of sex and
ethnicity. We have not tested either conjecture.

% ===========================================================================
\section{Who May Say What}
\label{sec:tiers}

\subsection{Issuer kinds in the signed list}

A trust list that admits an issuer usually admits its key for any statement. If
employers are to issue employment records, that is a problem: a key admitted so
that a small company can confirm its staff should not let it confirm a degree.
ACTIK signs each issuer's kind into the list and refuses any other credential
type from an employer (\kind{specified}).

\textbf{Proposition 5} (\kind{analytic}). \emph{If issuer kinds are recorded only
in a database, then whoever can write that database can promote an issuer; if
they are signed into the list, only the Root can.} It follows from the premise
that verifiers trust the Root's signature and not the database.

\subsection{Records people will carry}

An employment record is a statement about a person made by someone with power
over them. ACTIK's record cannot carry salary, reason for leaving, performance or
discipline, and the employee must accept it before it reaches their wallet
(\kind{specified}).

\textbf{Proposition 6} (\kind{conjectural}). \emph{Employees will accept and
present employment records more readily when the format guarantees that a record
cannot carry anything that could be used against them.} The cost is that a record
says less than a reference letter does. We think that is the right side to err
on, and we have not asked anyone.

\subsection{``Current'' has a date}

A record saying \emph{currently employed} is true on the day it is signed and
becomes false without anyone touching it. ACTIK never displays it without that
date (\kind{specified}), and asks employers to replace and withdraw a record when
the job ends. A verifier is told how fresh the withdrawal list it checked is, and
offline verifiers can be up to the list's refresh interval behind the issuer
(\kind{analytic}).

\subsection{Withdrawal and binding}

The withdrawal list names no one: each entry is a hash, so a reader learns
nothing unless they already hold the document number or credential identifier it
was computed from (\kind{carried}~\citep{chay2026qrseal}). Where document numbers
are sequential, a reader can hash candidates and learn \emph{that} a number was
withdrawn, never whose it was or why (\kind{analytic}). Binding removes two
attacks: presenting a copied credential, and replaying an answer made for one
requester to another (\kind{specified}, tested). It does not remove collusion: a
holder who shares their PIN, or hands over an unlocked phone, can have someone
else present for them (\kind{analytic}). That is why the name is still shown and
an identity check at interview is still advised.

% ===========================================================================
\section{Discussion}
\label{sec:discussion}

\subsection{What the paper adds}

The parts are known; the paper adds an arrangement and an account of it. A
verifier contract without a verdict, enforced by types and tests. A fixed limit
on what may be asked, written into the request format and enforced at three
points, suited to a setting without a register of requesters. Issuer kinds signed
into the trust list, and an employment record designed so that it cannot be used
against its holder. And a statement, for each, of which claims are true of the
code and which are guesses.

\subsection{Beside the national platform: more, and a choice}
\label{sec:choice}

ACTIK is not proposed against Verify.gov.kh. It is proposed as a second option
that does more of what a holder needs, for people and institutions to choose
between. Table~\ref{tab:choice} sets the two side by side. The rows for
Verify.gov.kh rest on its own published description~\citep{verifygovkh} and on
the earlier analysis~\citep{chay2026qrseal} (\kind{carried}); the rows for ACTIK
are \kind{specified}.

\begin{widetable}
\centering\small
\caption{What each option gives a holder and a verifier. Neither is required by the
other; an issuer may use both, and a holder may present whichever a verifier
accepts.}
\label{tab:choice}
\begin{tabularx}{\linewidth}{@{}lXX@{}}
\toprule
& Verify.gov.kh & ACTIK \\
\midrule
What the holder has & a document with a code that points to the platform's record & a signed credential in their own encrypted wallet, and optionally a printed code that carries the credential itself \\
Who can issue & recognised government bodies, as the platform describes itself & any issuer the Root admits: accredited institutions, including training providers, and registered employers for employment records \\
How it is checked & online, against the platform's servers & offline, against signed lists the verifier already holds \\
What is shown & the stored contents of the record & what the holder chooses to disclose, above a fixed minimum \\
Proof that the presenter is the holder & none in the code; a copy of the link works & bound credentials need the holder's key, for that verifier only \\
Withdrawal & visible at once to the next reader & visible as of the issuer's last signed list, up to its refresh interval later \\
Standing & a national government service in operation & a design with tests, not deployed \\
\bottomrule
\end{tabularx}
\end{widetable}

Two rows favour the national platform and we state them plainly. A withdrawal on
an online record takes effect at the next lookup, where ACTIK's takes effect only
when a verifier next refreshes the issuer's list. And Verify.gov.kh is an
operating government service with institutional standing; ACTIK is a tested
design that nobody yet uses.

The choice is not exclusive. A university may keep issuing through the national
platform and also sign an ACTIK credential for the same degree; a training
provider whose courses the national platform does not cover can issue through
ACTIK alone; a holder may present whichever the verifier in front of them
accepts. We propose no integration between the two and claim none.

\textbf{Proposition 7} (\kind{conjectural}). \emph{Given the choice, holders of
short courses and employment records, for which no national record exists, will
take up a holder-kept credential more readily than holders of degrees, which the
national platform already covers.} Nothing here tests it.

\subsection{The ecosystem}

Within CamboVerse, ACTIK is the layer that holds data and CamboVerse the layer
that shows it, joined by a portable file the holder exports rather than a live
interface. This keeps documents and personal data out of the display layer and
means that the display layer never needs to know who looks at whose exhibit. The
same arrangement --- a holder-controlled data layer, a display layer that verifies
what it is given, and a file between them --- is available to any platform that
wants to show people's achievements without holding them (\kind{analytic}).

% ===========================================================================
\section{What We Do Not Know}
\label{sec:unknown}

We would rather list the paper's limits than have them found.

There is no empirical content about people. Propositions~2, 4, 6 and~7 are
conjectures. Propositions~1, 3 and~5 follow from stated premises and are not
thereby true of the world.

There is no deployment. The database rules were tested against a local instance
with a stubbed authentication schema, and the interface against a stand-in for
the hosted database. Neither is evidence about the hosted service.

No measurement was taken on a low-cost phone. Table~\ref{tab:measure} reports a
server processor. The design targets phones on mobile networks, and a 1{,}000-entry
withdrawal list is about 118\,KB, which grows linearly with withdrawals; we have not
measured what that costs on the devices that matter.

The Root is a role nobody holds. Who should hold it, how institutions and
employers are admitted, and who republishes the list every month are governance
questions this paper does not answer. With many issuers, a verifier is in the end
trusting whoever admits them.

The labour-law list in Section~\ref{sec:known-hiring} is quoted from secondary
summaries. Holder binding covers only credentials issued after the holder had a
wallet; earlier ones stay unbound. And the Khmer text of the interface has not yet
been reviewed by native speakers.

% ===========================================================================
\section{What Should Happen Next}
\label{sec:next}

Four studies would move this materially, in the order we would run them.

First, a measurement on two or three low-cost Android phones on a mobile network:
time to open the trust list, verify a bound share, verify a printed code, and
fetch a withdrawal list of realistic size. It can be run now, needs no
participants, and would replace Table~\ref{tab:measure}'s server figures with the
ones that matter.

Second, a reading study of the verification screen: participants shown either
ACTIK's evidence layout or a single success mark, across a short course, a degree,
a withdrawn credential and an unchecked one, asked what they would conclude. It
tests Proposition~2 directly.

Third, structured interviews with a small number of employers and registrars about
the request format and employment records: what they would ask that the format
forbids, and whether employees would accept records under these rules. It bears on
Propositions~4 and~6 without claiming to test them.

Fourth, a pilot with one issuer and one employer who already hires its graduates,
once the governance of the Root role is settled. Only a pilot can say whether any
of this changes who gets hired.

% ===========================================================================
\section{Conclusion}
\label{sec:conclusion}

Someone who finishes a course should be able to keep the certificate, show it to
anyone, and show only what is needed. Cambodia's national platform verifies
official records online; ACTIK puts a signed credential in the holder's hand, and
is offered beside the platform for people to choose. A certificate check that
says \emph{verified} answers a narrower question than it
appears to, and a request for proof that can ask anything will be used to ask too
much. Both are settled by format before anyone looks at a credential. ACTIK's
answer is to return evidence instead of a verdict, to write the limits on asking
into the request itself, to sign into the trust list what kind of issuer each
issuer is, and to bind credentials to their holders where it can. None of the
mechanisms is new. Their arrangement is, and so is the discipline of saying, for
each claim, whether the code makes it true or whether it is a hope. The code is
public; the hopes are untested.

% ===========================================================================
\section*{Declarations}

\paragraph{Conflict of interest.} The authors declare no commercial or financial
relationships that could be construed as a potential conflict of interest. Both
authors hold appointments at the National University of Management. The first
author directs the CamboVerse Center, an institute at that University; the second
is an engineer on the CamboVerse project. These are affiliations. The University
issues degrees and uses the national verification platform discussed in
Section~\ref{sec:known}. Participation here is as individual researchers; the
paper is not University output, was not reviewed or approved under any University
process, and should not be cited as a University position.

\paragraph{Funding.} This research received no external funding.

\paragraph{Author contributions.} \emph{To be confirmed by the authors before
posting.} SC conceived the study and the claims, and drafted the manuscript. MV
built the ACTIK application on which Section~\ref{sec:system} reports, as recorded
in the repository history, and verified the description of it. Both authors
reviewed and approved the submitted version and accept responsibility for its
content.

\paragraph{Relationship to prior work.} The printed-certificate format and the
withdrawal-list entry design are carried from a preprint by the first
author~\citep{chay2026qrseal}; the method, including the four kinds of claim, from
a companion CamboVerse case by several of the present
authors~\citep{chay2026precinct}; and the observation about controls in layers
from a further companion paper~\citep{chay2026stages}. None of them is a source of
evidence for the claims made here.

\paragraph{Data and code availability.} No data about people were generated or
analysed. The source code, tests, database schema and measurement script are
published at \url{https://github.com/Actik-Cam/ACTIK}~\citep{actik2026repo}. Every
\kind{specified} claim in Section~\ref{sec:system} can be checked against it, the
counts in Table~\ref{tab:tests} are reproduced by \code{npm test}, and the figures
in Table~\ref{tab:measure} by \code{npm run measure}. Test fixtures use fictional
people; some fixtures and documentation name real institutions as illustrations
only, which implies no relationship with them.

\paragraph{Generative AI disclosure.} The authors used Anthropic's Claude during
this work: to implement parts of the system described in Section~\ref{sec:system}
(the trust layer, withdrawal lists, printed certificates, proof requests,
employment records, holder binding and their tests, as recorded in the repository
history), to search for prior art, to verify reference metadata, and to draft and
typeset this manuscript. The search for prior art withdrew three claims we had intended to make
(Section~\ref{sec:method}). The authors reviewed the AI-assisted
code and text, checked the cited sources, and take full responsibility for the
content of this paper.

\paragraph{Disclaimer.} This is a preprint and has not been peer reviewed.

\bibliographystyle{plainnat}
\bibliography{references}

\end{document}
